\documentclass[12pt]{article}

\usepackage[utf8]{inputenc}
\usepackage[T1]{fontenc}
\usepackage[portuguese]{babel}
\usepackage{geometry}
\geometry{a4paper, margin=2.5cm}

\usepackage{amsmath}
\usepackage{amssymb}
\usepackage{hyperref}
\usepackage{enumitem}
\usepackage{tcolorbox}

\hypersetup{
    colorlinks=true,
    linkcolor=blue,
    filecolor=magenta,      
    urlcolor=blue,
}

\begin{document}

\begin{center}
    \vspace{0.5cm}
    \textbf{INSTITUTO FEDERAL DO CEARÁ -- CAMPUS TAUÁ} \\
    \textbf{Tecnologia em Análise e Desenvolvimento de Sistemas} \\
    \textbf{Disciplina: Ciência de Dados / Análise de Dados} \\
    \textbf{Professor: Reginaldo Fernandes} \\
    \vspace{0.5cm}
    \large{\textbf{Atividade de Fixação 3 -- Testes de Permutação e Causalidade (N2)}}
    \vspace{0.3cm}
\end{center}

\hrule
\vspace{0.4cm}

\section*{Apresentação da Atividade}
Esta atividade de fixação aborda a lógica comparativa com dados categóricos via TVD, as bases do método científico (estudos observacionais, RCTs e resultados potenciais), e as discussões contemporâneas sobre replicação científica e uso de $p$-valores.

\vspace{0.3cm}
\hrule
\vspace{0.5cm}

\section*{Questões Teóricas e Computacionais}

\begin{enumerate}[label=\textbf{Questão \arabic*}]
    \item O caso \textit{Robert Swain vs. Alabama (1965)} e o julgamento em \textit{Alameda County} evidenciaram o papel da estatística na avaliação da representatividade racial em júris.
    \begin{enumerate}[label=\alph*)]
        \item Defina matematicamente a estatística de teste Total Variation Distance (TVD). Para que tipo de variáveis essa estatística é ideal?
        \item Suponha que a distribuição racial de um município cearense seja: $70\%$ Parda, $20\%$ Branca e $10\%$ Preta. Um júri sorteado de 100 pessoas apresentou a seguinte composição: $50\%$ Parda, $40\%$ Branca e $10\%$ Preta. Calcule o valor do TVD observado para esta amostra.
    \end{enumerate}

    \item Diferencie um \textbf{Estudo Observacional} de um \textbf{Experimento Controlado Aleatório (RCT)}. Por que apenas o segundo permite inferir causalidade direta de um tratamento?

    \item Explique o conceito de \textbf{Fator de Confundimento (Confounding Variable)} e dê um exemplo de como ele pode invalidar as conclusões de um estudo de correlação estatística em saúde pública.

    \item Explique o modelo de \textbf{Resultados Potenciais (Potential Outcomes)} através da metáfora didática do ``bilhete de dois lados''. \\
    Por que dizemos que o problema fundamental da inferência causal é a impossibilidade de ler as duas faces do bilhete para o mesmo paciente?

    \item A meta-análise é uma ferramenta fundamental no método científico contemporâneo.
    \begin{enumerate}[label=\alph*)]
        \item O que é uma Meta-Análise e por que ela oferece maior robustez científica do que estudos clínicos individuais com amostras pequenas?
        \item Relacione a meta-análise com o conceito de replicação científica.
    \end{enumerate}

    \item Muitas revistas científicas e sociedades de estatística têm criticado o uso excessivo e mecânico de $p$-valores na tomada de decisão (especialmente a barreira arbitrária de $p < 0.05$).
    \begin{enumerate}[label=\alph*)]
        \item Explique dois problemas práticos decorrentes do uso inadequado de $p$-valores (ex: \textit{p-hacking}, viés de publicação).
        \item Quais informações adicionais um analista de dados deve sempre reportar ao divulgar o resultado de um teste de hipóteses para garantir o rigor estatístico?
    \end{enumerate}
\end{enumerate}

\pagebreak

\section*{Gabarito Comentado e Dicas de Resolução}

\begin{tcolorbox}[colback=yellow!10!white,colframe=yellow!50!black,title=\textbf{Dica de Estudo}]
Use as resoluções comentadas abaixo para verificar suas respostas e orientar sua preparação para a avaliação teórica presencial.
\end{tcolorbox}

\begin{description}
    \item[Resolução 1:] 
    \begin{enumerate}[label=\alph*)]
        \item O TVD (Total Variation Distance) é a metade da soma das diferenças absolutas entre as proporções observadas e as esperadas de cada categoria: $TVD = \frac{1}{2}\sum |P_i^{\text{obs}} - P_i^{\text{esp}}|$. É ideal para comparar distribuições de variáveis categóricas.
        \item Cálculo do TVD:
        \begin{align*}
        TVD &= \frac{1}{2} \left( |0.50 - 0.70| + |0.40 - 0.20| + |0.10 - 0.10| \right) \\
        &= \frac{1}{2} (0.20 + 0.20 + 0.00) = \mathbf{0.20}
        \end{align*}
    \end{enumerate}

    \item[Resolução 2:] Em estudos observacionais, os indivíduos decidem sua exposição, o que permite a presença de variáveis de confundimento que distorcem o efeito causal. Em Experimentos Controlados Aleatórios (RCTs), a alocação aleatória sorteia a exposição de cada paciente, distribuindo uniformemente as variáveis ocultas entre controle e tratamento, neutralizando o confundimento e isolando o efeito direto do tratamento.

    \item[Resolução 3:] Fator de confundimento é uma variável não observada que está correlacionada tanto com o tratamento quanto com o resultado de interesse. Exemplo clássico: Observa-se que pessoas que consomem mais café apresentam maior incidência de problemas cardíacos. Porém, o tabagismo atua como fator de confundimento (fumantes consomem mais café e o cigarro causa os problemas cardíacos). Se não controlarmos o cigarro, culparemos erroneamente o café.

    \item[Resolução 4:] Pelo modelo de resultados potenciais, cada participante possui dois resultados possíveis: o resultado sob tratamento e sob controle (bilhete de dois lados). O problema fundamental da inferência causal é que, ao alocarmos o paciente em um grupo, só podemos ver um dos lados da resposta (a outra face torna-se contrafactual/desconhecida). Por isso, inferir causalidade a nível individual é impossível, restando-nos comparar as médias agregadas de grupos estatisticamente idênticos.

    \item[Resolução 5:] 
    \begin{enumerate}[label=\alph*)]
        \item Meta-análise é uma metodologia estatística para combinar os resultados quantitativos de múltiplos estudos independentes. Ela aumenta o tamanho amostral acumulado ($N$), o que eleva consideravelmente o poder do teste estatístico e reduz a margem de erro.
        \item A meta-análise valida a replicação: ela verifica se diferentes experimentos em contextos distintos confirmam sistematicamente o mesmo efeito biológico ou clínico.
    \end{enumerate}

    \item[Resolução 6:] 
    \begin{enumerate}[label=\alph*)]
        \item O \textit{p-hacking} consiste na realização de múltiplos testes sucessivos, alteração de critérios de exclusão de dados ou inclusão de variáveis até que um $p$-valor ficticiamente significativo ($p < 0.05$) seja encontrado. O \textit{viés de publicação} ocorre porque periódicos científicos priorizam a publicação de estudos com resultados positivos significativos, gerando distorções na literatura médica (estudos com efeitos nulos são engavetados).
        \item O analista de dados deve sempre reportar: (1) O P-valor exato, (2) O tamanho da amostra ($N$), (3) O tamanho do efeito (effect size) e (4) O método estatístico com todas as suas premissas declaradas.
    \end{enumerate}
\end{description}

\end{document}
